\section{Limitations of Vector-Based Retrieval}
\label{sec:gr-vector-limitations}

Although dense vector retrieval is effective for locating passages that are
semantically similar to a query, it treats the corpus as a flat collection of
independent chunks and discards the relational structure that connects
them~\cite{peng2024graphrag}. This flat view is the root of several documented
failure modes, and it is the primary motivation for introducing graph structure
into retrieval~\cite{peng2024graphrag, edge2024graphrag}.

\subsection{Multi-Hop and Relational Queries}

Questions whose answers require chaining several facts together are difficult for
naive RAG, because the individual passages that must be combined are often not
individually similar to the query and may not co-occur in any single
chunk~\cite{peng2024graphrag}. Multi-hop question answering benchmarks were
created precisely to expose this weakness, requiring a system to find and reason
over multiple supporting documents rather than retrieving a single relevant
passage~\cite{yang2018hotpotqa}. Because vector similarity captures topical
closeness but not the explicit relations between entities, it provides no
reliable way to follow such reasoning chains~\cite{peng2024graphrag}.

\subsection{Global vs. Local Questions}

A second limitation concerns \emph{global} or thematic questions directed at an
entire corpus, such as ``what are the main themes in this dataset?''.
Such questions are inherently query-focused summarization tasks rather than
retrieval tasks, and top-$k$ passage retrieval cannot answer them because no
small set of chunks contains the global answer~\cite{edge2024graphrag}. Naive RAG
consequently underperforms on comprehensiveness and diversity for this class of
question, whereas a graph-structured index that pre-computes summaries of related
entities was shown to improve substantially on both
axes~\cite{edge2024graphrag}.
